problem:  
Let \( S \) be a solvable Lie group with Lie algebra \( \mathfrak{s} = \mathfrak{a}\oplus\mathfrak{n} \), such that the adjoint action of \( \mathfrak{a} \) on \( \mathfrak{n} \) is **not of real type**—that is, some \( \operatorname{ad}_A|_{\mathfrak{n}} \) have non‑real complex eigenvalues.  
Denote by \( [\cdot,\cdot]_t \) the solution to the **normalized bracket flow**
\[
\frac{d}{dt}[\cdot,\cdot]_t = -\pi(Q([\cdot,\cdot]_t))([\cdot,\cdot]_t),
\]
acting on the variety of solvable Lie brackets on the fixed vector space \( \mathfrak{s} \), where \( Q([\cdot,\cdot]) \) is the Ricci operator and \( \pi \) the natural representation of \( \mathrm{End}(\mathfrak{s}) \) on brackets.  
Equilibria of this system are solvsoliton brackets, which correspond to Ricci soliton metrics on \( S \).

> **Conjecture (Existence of Non‑Convergent Homogeneous Ricci‑Flow Trajectories in the Non‑Real‑Type Regime).**  
> There exists a solvable Lie algebra \( \mathfrak{s} \) not of real type such that the normalized bracket‑flow trajectory \( [\cdot,\cdot]_t \) starting from some initial bracket does **not** converge (modulo automorphisms and scaling) to a single solvsoliton.  
>  
> Equivalently, non‑real‑type solvmanifolds admit homogeneous Ricci‑flow evolutions whose asymptotic limit set is not a single orbit of the solvsoliton type—potentially exhibiting recurrent, oscillatory, or multiple‑attractor behavior.

Why is it a "good" problem:  
1. **Mathematical soundness and precision:**  
   The problem is logically well‑formed: non‑real‑type solvable Lie algebras are known to exist, the normalized bracket flow is a canonical system encoding homogeneous Ricci flow, and solvsolitons are its equilibria. The question isolates a sharp, verifiable phenomenon—*non‑convergent long‑time dynamics*—within a well‑defined finite‑dimensional differential system.

2. **Genuine open frontier:**  
   For solvable Lie algebras of real type, convergence of the normalized bracket flow to a unique solvsoliton is established. Beyond this restrictive class, no general asymptotic theorem is known. Whether convergence fails in the non‑real‑type case is a **major unresolved gap** in the current Ricci‑flow theory on homogeneous spaces.

3. **Bridge across disciplines:**  
   The conjecture links **geometric analysis** (Ricci flow), **algebraic geometry** (variety of Lie brackets), and **nonlinear dynamics** (limit sets and recurrence). Establishing a counterexample would introduce chaotic‑like or multi‑attractor phenomena into geometric flows, which historically behave as gradient systems—a paradigm shift uniting geometry and dynamical systems.

4. **Profound simplicity and reach:**  
   The conjecture is compact (“Does the flow always converge?”) but demands deep insight into how complex eigenstructures of \( \operatorname{ad}_{\mathfrak{a}} \) sabotage gradient‑like properties in the bracket flow. Whether such algebraic complexity can translate into analytic non‑convergence is a profound, conceptually central question.

5. **Potential impact:**  
   - A **positive answer** (existence of non‑convergent trajectories) would reveal fundamentally new dynamical behaviors in Ricci flow, challenging the prevailing convergence paradigm for homogeneous geometries.  
   - A **negative answer** (universal convergence) would confirm that solvsolitons are global attractors even beyond real type, thus completing the characterization of homogeneous Ricci‑flow asymptotics.

This problem epitomizes the best traits of a high‑level mathematical challenge: it is concise, structurally transparent, deeply interdisciplinary, and located exactly at a critical boundary of current knowledge, making it an outstanding “good” research problem.